\documentclass[11pt]{article}

\usepackage[T1]{fontenc}
\usepackage[utf8]{inputenc}
\usepackage{lmodern}
\usepackage[margin=1.03in]{geometry}
\usepackage{microtype}
\usepackage{amsmath,amssymb,amsthm,mathtools}
\usepackage{booktabs,longtable,tabularx,array}
\usepackage{enumitem}
\usepackage{xcolor}
\usepackage{hyperref}
\usepackage[nameinlink,capitalise,noabbrev]{cleveref}

\hypersetup{
  colorlinks=true,
  linkcolor=blue!45!black,
  citecolor=green!35!black,
  urlcolor=blue!55!black,
  pdftitle={Moving-Prime Arithmetic and a Kummer-Period Bridge for the Alaoglu--Erdos Problem},
  pdfauthor={OpenAI GPT-5.6 Pro}
}

\allowdisplaybreaks
\setlist[itemize]{topsep=4pt,itemsep=2pt,parsep=1pt}
\setlist[enumerate]{topsep=4pt,itemsep=2pt,parsep=1pt}
\renewcommand{\arraystretch}{1.12}

\newtheorem{theorem}{Theorem}[section]
\newtheorem{proposition}[theorem]{Proposition}
\newtheorem{lemma}[theorem]{Lemma}
\newtheorem{corollary}[theorem]{Corollary}
\newtheorem{criterion}[theorem]{Criterion}
\newtheorem{conjecture}[theorem]{Conjecture}
\theoremstyle{definition}
\newtheorem{definition}[theorem]{Definition}
\newtheorem{problem}[theorem]{Problem}
\newtheorem{program}[theorem]{Program}
\theoremstyle{remark}
\newtheorem{remark}[theorem]{Remark}
\newtheorem{warning}[theorem]{Warning}

\newcommand{\N}{\mathbb N}
\newcommand{\Z}{\mathbb Z}
\newcommand{\Q}{\mathbb Q}
\newcommand{\R}{\mathbb R}
\newcommand{\Gm}{\mathbb G_m}
\newcommand{\vp}{v_p}
\newcommand{\ord}{\operatorname{ord}}
\newcommand{\lcm}{\operatorname{lcm}}
\newcommand{\Sym}{\operatorname{Sym}}
\newcommand{\Div}{\operatorname{Div}}
\newcommand{\Int}{\operatorname{Int}}
\newcommand{\per}{\operatorname{per}}
\newcommand{\Span}{\operatorname{span}}
\newcommand{\supp}{\operatorname{supp}}
\newcommand{\qbinom}[3]{\genfrac{[}{]}{0pt}{}{#1}{#2}_{#3}}
\newcommand{\AE}{Alaoglu--Erd\H{o}s}
\newcommand{\e}{\mathrm e}
\newcommand{\dd}{\,\mathrm d}

\title{\textbf{Moving-Prime Arithmetic and a Kummer-Period Bridge\\
for the Alaoglu--Erd\H{o}s Problem}\\[0.7em]
\large A nonduplicative progress report with exact geometric $p$-orderings,\\
a sharp half-mass barrier, a gcd-compatibility endgame, and a motivic conditional proof}
\author{OpenAI GPT-5.6 Pro}
\date{22 August 2026}

\begin{document}
\maketitle

\begin{abstract}
Let $x\in\R$ and suppose that $M=2^x$ and $N=3^x$ are integers.  The
\AE{} conjecture asserts that $x$ is an integer.  This report does not claim
an unconditional proof.  It records a set of new results and proof targets
that were selected specifically to avoid reproducing the broad determinant,
Mahler-function, fixed-prime $p$-adic, interpolation, approximation-theoretic,
and model-theoretic developments of the supplied Report~A.

The main unconditional result is an exact simultaneous $p$-ordering theorem
for every integral geometric progression $\{a q^n:n\geq0\}$.  The natural
ordering works for every prime at once, and its generalized factorial is
\[
 a^kq^{\binom{k}{2}}\prod_{r=1}^k(q^r-1).
\]
This gives a canonical Newton--Gaussian basis for all rational polynomials
integer-valued on the progression.  An exact local valuation formula then
splits the logarithmic denominator mass into two asymptotically equal parts:
one half is concentrated at the fixed primes dividing $q$, while the other
half is carried by the cyclotomic factors $q^r-1$.  On the full geometric
Vandermonde scale the same phenomenon occurs at cubic order.  Moreover, the
cyclotomic half escapes every fixed finite set of primes.  Consequently, a
fixed-$2$/$3$-adic determinant argument can capture at most one half of the
universal geometric denominator at leading order; the missing arithmetic
must move with the interpolation scale.

To turn this diagnosis into an endgame, the report introduces a weighted
moving-prime compatibility mass
\[
 \mathfrak C_d(q,Q)=\sum_{r=1}^d(d+1-r)
        \log\gcd(q^r-1,Q^r-1).
\]
A theorem of Bugeaud--Corvaja--Zannier implies a sharp zero--one law when $q$
is prime: $\mathfrak C_d(q,Q)$ has full cubic mass if $Q$ is an integral
power of $q$, and is $o(d^3)$ if $q$ and $Q$ are multiplicatively
independent.  Therefore the original conjecture follows from one concrete
bridge statement: a common real exponent for $(2,M)$ and $(3,N)$ must force a
positive proportion of moving-prime compatibility in at least one pair.
A block-Vandermonde lemma shows exactly how such gcd mass can be amplified in
integer determinants.

A second, logically independent route identifies the sparse logarithmic
determinant with the numerical period of a symmetric tensor of Kummer
logarithms.  Its weight-$(1,1)$ motivic coaction is precisely the formal prime
symbol whose vanishing forces $M=2^k$ and $N=3^k$.  This yields a clean
conditional proof under a restricted weight-two period-injectivity statement,
and explains why the symmetric square, rather than the raw tensor square, is
the correct receptacle.  A function-field divisor theorem supplies an
unconditional analogue and isolates the constant-field obstruction to
specialization.  The final sections give a focused experimental agenda and a
Lean formalization plan for every unconditional component.
\end{abstract}

\tableofcontents
\newpage

\section{Scope, status, and the nonduplication boundary}

The problem is
\begin{equation}\label{eq:AE}
 2^x,3^x\in\Z \qquad\Longrightarrow\qquad x\in\Z,
 \qquad x\in\R.
\end{equation}
It originated in the work of Alaoglu and Erd\H{o}s on highly composite and
related integers \cite{AlaogluErdos1944}.  In exponential form, a putative
counterexample supplies the algebraic $2\times2$ array
\[
 \exp\!\begin{pmatrix}
     \log 2 & \log 3\\
     x\log 2 & x\log 3
 \end{pmatrix}
 =\begin{pmatrix}2&3\\M&N\end{pmatrix},
 \qquad M,N\in\Z_{>0},
\]
whose logarithmic matrix has rank one.  This is a highly structured positive
integral instance of the four-exponentials problem.

The supplied Report~A already explores a large collection of determinant,
fixed-prime $p$-adic, Mahler, interpolation, approximation, and logical
routes.  Its most useful common reduction is the free-monoid amplification:
if $M=2^x$ and $N=3^x$, then
\begin{equation}\label{eq:monoid-map}
  (2^a3^b)^x=M^aN^b\in\Z_{>0}
  \qquad(a,b\in\N_0).
\end{equation}
The present report imports only this compact baseline and does not repeat the
prior surveys or determinant calculations.  Its new project-level content is
concentrated in four directions:
\begin{enumerate}
 \item exact simultaneous $p$-orderings and generalized factorials of
       geometric progressions;
 \item a rigorous half-mass and escape-of-mass theorem, including its full
       Vandermonde version;
 \item a moving-prime gcd compatibility functional with a zero--one theorem
       and an explicit determinant amplifier;
 \item a Kummer-period/coaction interpretation and a function-field divisor
       analogue.
\end{enumerate}

The competitive claims mentioned in the task are not used as mathematical
inputs.  Every assertion below is either proved in the report, explicitly
marked conditional, or attached to a primary mathematical reference.

\subsection{Immediate reductions}

Put
\[
 M=2^x,\qquad N=3^x.
\]
If $x<0$, then $0<2^x<1$, so $2^x$ cannot be a positive integer.  Thus
$x\geq0$.  The case $x=0$ is trivial.

\begin{lemma}[Rational and algebraic cases]\label{lem:rational-algebraic}
If $2^x\in\Z$ and $x\in\Q$, then $x\in\Z$.  Consequently, a nonintegral
counterexample to \eqref{eq:AE} must be transcendental.
\end{lemma}

\begin{proof}
Write $x=a/b$ with $a,b\in\N$ and $\gcd(a,b)=1$.  If $M=2^{a/b}\in\Z$, then
$M^b=2^a$.  Taking the $2$-adic valuation gives $b\,v_2(M)=a$, hence $b\mid a$
and $b=1$.  If $x$ were algebraic irrational, the Gelfond--Schneider theorem
would make $2^x$ transcendental; see, for example,
\cite[Ch.~1]{Waldschmidt2000}.  Therefore a nonintegral counterexample is
transcendental.
\end{proof}

The exact real relation is the logarithmic determinant
\begin{equation}\label{eq:logdet}
 \mathcal D(M,N):=\log 2\,\log N-\log 3\,\log M=0.
\end{equation}
The report will attack \eqref{eq:logdet} in two different ways.  The
moving-prime route asks how much congruence structure is preserved by the
monoid map in \eqref{eq:monoid-map}.  The period route asks whether the
numerical vanishing in \eqref{eq:logdet} can be lifted to a formal relation
between prime logarithm symbols.

\section{Exact simultaneous $p$-orderings of geometric progressions}

\subsection{Definitions}

Let $S\subset\Z$ and let $p$ be prime.  A sequence
$s_0,s_1,s_2,\ldots$ in $S$ is a \emph{$p$-ordering} if, at stage $k$, the
quantity
\[
 v_p\!\left(\prod_{j=0}^{k-1}(s_k-s_j)\right)
\]
is minimal among the corresponding quantities obtained by replacing $s_k$
with an unused element of $S$.  Bhargava introduced this construction and
its associated generalized factorials in a much broader Dedekind-domain
setting \cite{Bhargava1997,Bhargava2000}.  A sequence that is a $p$-ordering
for every prime simultaneously will be called a \emph{simultaneous
$p$-ordering}.

For integers $a,q\geq1$ with $q\geq2$, write
\[
 G(a,q):=\{a q^n:n\in\N_0\}.
\]
The key observation is that the ordinary exponent ordering of $G(a,q)$ is
simultaneously optimal at every prime.  The proof is an exact Gaussian
binomial identity, not an asymptotic estimate.

\subsection{The Gaussian divisibility identity}

For $0\leq k\leq n$, define the Gaussian binomial coefficient by
\begin{equation}\label{eq:qbinom}
 \qbinom{n}{k}{q}
 :=\prod_{r=1}^{k}\frac{q^{n-k+r}-1}{q^r-1}.
\end{equation}
It is an integer.  One proof uses the recurrence
\[
 \qbinom{n}{k}{q}
 =\qbinom{n-1}{k}{q}+q^{n-k}\qbinom{n-1}{k-1}{q}
\]
with the usual boundary conditions; see \cite[Ch.~3]{Andrews1976}.

\begin{theorem}[Exact simultaneous ordering]\label{thm:simultaneous-ordering}
Let $a,q\in\Z_{>0}$ and $q\geq2$.  The sequence
\[
 a,aq,aq^2,aq^3,\ldots
\]
is a $p$-ordering of $G(a,q)$ for every prime $p$.  More precisely, for all
$n\geq k$,
\begin{equation}\label{eq:ratio-qbinom}
 \frac{\displaystyle\prod_{j=0}^{k-1}(aq^n-aq^j)}
      {\displaystyle\prod_{j=0}^{k-1}(aq^k-aq^j)}
 =\qbinom{n}{k}{q}\in\Z.
\end{equation}
Consequently the $k$th generalized factorial is represented by
\begin{equation}\label{eq:geom-factorial}
 F_k(a,q):=\prod_{j=0}^{k-1}(aq^k-aq^j)
 =a^kq^{\binom{k}{2}}\prod_{r=1}^{k}(q^r-1),
 \qquad F_0(a,q):=1.
\end{equation}
\end{theorem}

\begin{proof}
For $n\geq k$,
\[
 \prod_{j=0}^{k-1}(aq^n-aq^j)
 =a^kq^{0+1+\cdots+(k-1)}
   \prod_{j=0}^{k-1}(q^{n-j}-1).
\]
The corresponding expression with $n=k$ has the same powers of $a$ and $q$.
Their quotient is
\[
 \frac{\prod_{r=n-k+1}^{n}(q^r-1)}
      {\prod_{r=1}^{k}(q^r-1)}
 =\qbinom{n}{k}{q},
\]
which is an integer.  Hence, for every prime $p$, the $p$-adic valuation at
the candidate $aq^n$ is at least the valuation at $aq^k$.  Candidates with
$n<k$ have already occurred and give a zero factor.  This proves simultaneous
minimality.  Formula \eqref{eq:geom-factorial} follows by setting $n=k$ and
reversing the product $k-j$.
\end{proof}

\begin{remark}[Why simultaneity matters]
A generic subset of $\Z$ has different greedy orderings at different primes.
Here one ordering is globally compatible because every comparison quotient is
an actual integer.  This converts local $p$-adic information into a single
integral denominator and permits exact global mass accounting.
\end{remark}

\subsection{A canonical Newton--Gaussian basis}

Define
\begin{equation}\label{eq:Bk}
 B_k^{(a,q)}(T)
 :=\frac{\prod_{j=0}^{k-1}(T-aq^j)}{F_k(a,q)}.
\end{equation}

\begin{theorem}[Regular basis for the integer-valued polynomial ring]
\label{thm:regular-basis}
For every $n,k\in\N_0$,
\begin{equation}\label{eq:B-evaluation}
 B_k^{(a,q)}(aq^n)=
 \begin{cases}
   0,&n<k,\\[2mm]
   \qbinom{n}{k}{q},&n\geq k.
 \end{cases}
\end{equation}
In particular $B_k^{(a,q)}$ is integer-valued on $G(a,q)$ and
$B_k^{(a,q)}(aq^k)=1$.  Every polynomial $P\in\Q[T]$ of degree at most $d$
that is integer-valued on $G(a,q)$ has a unique expansion
\begin{equation}\label{eq:regular-expansion}
 P(T)=\sum_{k=0}^{d}c_k B_k^{(a,q)}(T),
 \qquad c_k\in\Z.
\end{equation}
\end{theorem}

\begin{proof}
Equation \eqref{eq:B-evaluation} is exactly
\eqref{eq:ratio-qbinom}, with the zero case coming from a repeated node.  For
the expansion, determine $c_k$ recursively from the values at
$a,aq,\ldots,aq^d$.  The evaluation matrix
$\bigl(B_k^{(a,q)}(aq^n)\bigr)_{0\leq n,k\leq d}$ is lower triangular with
ones on the diagonal and integral entries.  Therefore the recursive
coefficients are integers.  The resulting polynomial agrees with $P$ at
$d+1$ distinct points, hence is equal to $P$.  Uniqueness follows from the
triangular matrix or from distinct degrees.
\end{proof}

This theorem supplies an exact denominator lattice for interpolation on a
geometric progression.  It is particularly relevant to \eqref{eq:monoid-map}:
every ray
\[
 3^b,\ 3^b2,\ 3^b2^2,\ldots
\]
and every ray
\[
 2^a,\ 2^a3,\ 2^a3^2,\ldots
\]
inside the $2$--$3$ smooth monoid has a canonical basis whose values and
denominators are known exactly.

\subsection{Exact Newton coefficients of a real power}

The preceding basis also gives a useful identity for the nonpolynomial
function $T\mapsto T^x$.

\begin{proposition}[Real-power divided difference]\label{prop:real-power-dd}
Let $q>1$, $x\in\R$, and put $Q=q^x$.  For every $k\geq0$,
\begin{equation}\label{eq:dd-power}
 [1,q,\ldots,q^k]T^x
 =\frac{\prod_{r=0}^{k-1}(Q-q^r)}{F_k(1,q)}.
\end{equation}
More generally,
\begin{equation}\label{eq:dd-power-scaled}
 F_k(a,q)\,[a,aq,\ldots,aq^k]T^x
 =a^x\prod_{r=0}^{k-1}(Q-q^r).
\end{equation}
Here brackets denote ordinary divided differences.
\end{proposition}

\begin{proof}
As a function of the auxiliary variable $Y=Q$, the divided difference
\[
 C_k(Y):=\sum_{j=0}^{k}\frac{Y^j}
 {\prod_{0\leq i\leq k,\ i\neq j}(q^j-q^i)}
\]
is a polynomial of degree $k$.  For $Y=q^r$ with $0\leq r<k$, the values
$Y^j=(q^j)^r$ come from a polynomial of degree $r$, so its $k$th divided
difference is zero.  Thus $C_k$ has roots $1,q,\ldots,q^{k-1}$.  Its leading
coefficient comes only from the $j=k$ term and equals $1/F_k(1,q)$.  This
proves \eqref{eq:dd-power}.  Scaling every node by $a$ multiplies a $k$th
divided difference of the homogeneous function $T^x$ by $a^{x-k}$; using
$F_k(a,q)=a^kF_k(1,q)$ gives \eqref{eq:dd-power-scaled}.
\end{proof}

For a putative \AE{} counterexample, \eqref{eq:dd-power-scaled} specializes
to the two separated families
\begin{align}
 F_k(3^b,2)\,[3^b,3^b2,\ldots,3^b2^k]T^x
   &=N^b\prod_{r=0}^{k-1}(M-2^r),\label{eq:ray2}\\
 F_k(2^a,3)\,[2^a,2^a3,\ldots,2^a3^k]T^x
   &=M^a\prod_{r=0}^{k-1}(N-3^r).\label{eq:ray3}
\end{align}
The right sides are integers.  These identities do not by themselves prove
that a divided difference is integral; instead, they show exactly where its
denominator resides and how the transverse smooth parameter separates from
the longitudinal geometric parameter.

\begin{proposition}[Fixed-base denominator dichotomy]
\label{prop:base-dichotomy}
Let $p\mid q$, put $h=v_p(q)$, and suppose that $Q=q^x\in\Z_{>0}$ is not an
integral power of $q$.  If $s=v_p(Q)$, then
\begin{equation}\label{eq:numerator-linear-vp}
 v_p\!\left(\prod_{r=0}^{k-1}(Q-q^r)\right)=sk+O_{p,q,Q}(1),
\end{equation}
and hence
\begin{equation}\label{eq:dd-quadratic-denom}
 v_p\bigl([1,q,\ldots,q^k]T^x\bigr)
 =-h\binom{k}{2}+sk+O_{p,q,Q}(1).
\end{equation}
\end{proposition}

\begin{proof}
If $rh>s$, then
\[
 v_p(Q-q^r)=v_p(Q)=s.
\]
If $rh<s$, then $v_p(Q-q^r)=rh$.  There is at most one index with $rh=s$;
its valuation is finite because $Q\neq q^r$.  Thus all but finitely many
terms have valuation $s$, proving \eqref{eq:numerator-linear-vp}.  Since
$p\mid q$ implies $p\nmid q^r-1$, formula \eqref{eq:geom-factorial} gives
$v_p(F_k(1,q))=h\binom{k}{2}$.  Subtracting valuations in
\eqref{eq:dd-power} proves \eqref{eq:dd-quadratic-denom}.
\end{proof}

The proposition is a rigorous warning: unless $Q$ is already a $q$-power,
the normalized divided differences acquire a quadratic denominator at the
base primes, while their numerator can cancel only linearly.  Any successful
argument must therefore use arithmetic beyond the fixed primes dividing the
base.

\section{Exact local valuation profiles}

The factorization in \eqref{eq:geom-factorial} permits a complete local
calculation.  This section records it because it distinguishes the fixed-base
mass from the moving cyclotomic mass with no hidden constants.

\begin{theorem}[Local valuation formula]\label{thm:local-profile}
Let $a,q\in\Z_{>0}$, $q\geq2$, and $k\geq1$.
\begin{enumerate}
 \item If $p\mid q$, then
 \begin{equation}\label{eq:local-pdivq}
   v_p(F_k(a,q))=k v_p(a)+\binom{k}{2}v_p(q).
 \end{equation}
 \item Let $p$ be odd and $p\nmid q$.  Put
 \[
   t=\ord_p(q),\qquad m=\left\lfloor\frac{k}{t}\right\rfloor,
   \qquad c=v_p(q^t-1).
 \]
 Then
 \begin{equation}\label{eq:local-podd}
   v_p(F_k(a,q))=k v_p(a)+mc+v_p(m!).
 \end{equation}
 \item If $p=2$ and $q$ is odd, put
 \[
   A=v_2(q-1),\qquad B=v_2(q+1),\qquad m=\lfloor k/2\rfloor.
 \]
 Then
 \begin{equation}\label{eq:local-two}
   v_2(F_k(a,q))=k v_2(a)+kA+mB+v_2(m!).
 \end{equation}
\end{enumerate}
\end{theorem}

\begin{proof}
The first assertion follows because $q^r-1$ is a $p$-adic unit when $p\mid q$.
For odd $p\nmid q$, the prime $p$ divides $q^r-1$ exactly when $t\mid r$.
Writing $r=tu$, the lifting-the-exponent lemma gives
\[
 v_p(q^{tu}-1)=v_p(q^t-1)+v_p(u)=c+v_p(u).
\]
Summing over $1\leq u\leq m$ yields \eqref{eq:local-podd}.  If $q$ is odd,
the $2$-adic lifting formula gives
\[
 v_2(q^r-1)=
 \begin{cases}
  A,&r\text{ odd},\\
  A+B+v_2(r)-1,&r\text{ even}.
 \end{cases}
\]
For $r=2u$ the second line is $A+B+v_2(u)$.  Summing the odd and even indices
produces $kA+mB+v_2(m!)$, proving \eqref{eq:local-two}.
\end{proof}

\begin{corollary}[Order-layer decomposition]\label{cor:order-layer}
For odd primes $p\nmid q$, all $p$-adic contribution to $F_k(1,q)$ is
controlled by the residue order $t=\ord_p(q)$ and the single initial depth
$v_p(q^t-1)$.  In particular, for fixed $p$,
\[
 v_p(F_k(1,q))=O_{p,q}(k).
\]
For every fixed finite set of primes $S$,
\begin{equation}\label{eq:fixed-cyc-linear}
 \sum_{\substack{p\in S\\p\nmid q}}v_p(F_k(1,q))\log p=O_{q,S}(k).
\end{equation}
\end{corollary}

The exact order-layer formula will reappear in the full Vandermonde setting,
where the weight $k+1-r$ raises the scale from quadratic to cubic.

\section{The half-mass law and escape from every fixed prime set}

\subsection{One-step generalized factorials}

Set
\[
 \Phi_k(q):=\prod_{r=1}^k(q^r-1),
 \qquad
 C_q:=\sum_{r=1}^{\infty}\log(1-q^{-r}).
\]
The series defining $C_q$ converges absolutely.

\begin{theorem}[Quadratic half-mass law]\label{thm:half-mass-factorial}
For fixed $q\geq2$,
\begin{align}
 \log F_k(1,q)
   &=k^2\log q+C_q+O_q(q^{-k}),\label{eq:total-factorial-mass}\\
 \sum_{p\mid q}v_p(F_k(1,q))\log p
   &=\binom{k}{2}\log q,\label{eq:base-factorial-mass}\\
 \log\Phi_k(q)
   &=\binom{k+1}{2}\log q+C_q+O_q(q^{-k}).\label{eq:cyc-factorial-mass}
\end{align}
Thus the primes dividing $q$ carry asymptotically one half of the logarithmic
mass of $F_k(1,q)$, and the cyclotomic factors $q^r-1$ carry the other half.
For fixed $a$, the extra term $k\log a$ is lower order and does not alter the
limit.
\end{theorem}

\begin{proof}
From \eqref{eq:geom-factorial},
\begin{align*}
 \log F_k(1,q)
 &=\binom{k}{2}\log q+
   \sum_{r=1}^k\bigl(r\log q+\log(1-q^{-r})\bigr)\\
 &=k^2\log q+\sum_{r=1}^k\log(1-q^{-r}).
\end{align*}
The tail of the convergent series is $O_q(q^{-k})$.  The base-prime identity
follows from unique factorization of $q$, and the cyclotomic identity is the
remaining factor.
\end{proof}

The asymptotic equality of the two halves might suggest that one can recover
the cyclotomic half from a few favorable small primes.  The next theorem
shows that this is impossible.

\begin{theorem}[Escape of cyclotomic mass]\label{thm:escape-factorial}
Fix $q\geq2$ and a finite set of primes $S$.  Then
\begin{equation}\label{eq:escape-factorial}
 \sum_{\substack{p\notin S\\p\nmid q}}
   v_p(F_k(1,q))\log p
 =\binom{k+1}{2}\log q+O_{q,S}(k).
\end{equation}
Equivalently, all but $O_{q,S}(k)$ of the quadratic cyclotomic mass is carried
by primes outside $S$.
\end{theorem}

\begin{proof}
The total cyclotomic mass is \eqref{eq:cyc-factorial-mass}.  The contribution
from the fixed set $S$ is $O_{q,S}(k)$ by
\cref{cor:order-layer}.  Subtraction proves the result.
\end{proof}

\begin{corollary}[No finite-prime recovery of the missing half]
\label{cor:no-finite-prime-factorial}
If $S$ contains every prime dividing $q$, then
\[
 \sum_{p\in S}v_p(F_k(1,q))\log p
 =\frac12\log F_k(1,q)+o(k^2).
\]
No enlargement of $S$ by finitely many primes changes the leading constant
$1/2$.
\end{corollary}

\subsection{Full geometric Vandermonde determinants}

For the nodes $a,aq,\ldots,aq^d$, define
\[
 V_d(a,q):=\prod_{0\leq i<j\leq d}(aq^j-aq^i).
\]

\begin{theorem}[Exact geometric Vandermonde factorization]
\label{thm:V-factorization}
For $d\geq1$,
\begin{equation}\label{eq:V-factorization}
 V_d(a,q)
 =a^{\binom{d+1}{2}}
  q^{\binom{d+1}{3}}
  \prod_{r=1}^{d}(q^r-1)^{d+1-r}.
\end{equation}
Moreover,
\begin{equation}\label{eq:V-product-factorials}
 V_d(a,q)=\prod_{k=1}^{d}F_k(a,q).
\end{equation}
\end{theorem}

\begin{proof}
Each difference is $aq^i(q^{j-i}-1)$.  There are
$\binom{d+1}{2}$ factors of $a$, and
\[
 \sum_{0\leq i<j\leq d}i
 =\sum_{i=0}^{d-1}i(d-i)=\binom{d+1}{3}.
\]
A difference $r=j-i$ occurs $d+1-r$ times, proving
\eqref{eq:V-factorization}.  Grouping the pairs by their larger index $j=k$
gives \eqref{eq:V-product-factorials}.
\end{proof}

\begin{theorem}[Cubic half-mass law]\label{thm:half-mass-V}
For fixed $q\geq2$ and $a=1$,
\begin{align}
 \log V_d(1,q)
 &=\frac{d(d+1)(2d+1)}6\log q+O_q(d),
 \label{eq:V-total-mass}\\
 \sum_{p\mid q}v_p(V_d(1,q))\log p
 &=\binom{d+1}{3}\log q,\label{eq:V-base-mass}\\
 \sum_{r=1}^{d}(d+1-r)\log(q^r-1)
 &=\binom{d+2}{3}\log q+O_q(d).
 \label{eq:V-cyc-mass}
\end{align}
The fixed-base and cyclotomic terms are both asymptotic to
$\frac16d^3\log q$, while the total is asymptotic to
$\frac13d^3\log q$.
\end{theorem}

\begin{proof}
Only the error term needs comment.  We have
\[
 \sum_{r=1}^{d}(d+1-r)\log(q^r-1)
 =\log q\sum_{r=1}^{d}(d+1-r)r
  +\sum_{r=1}^{d}(d+1-r)\log(1-q^{-r}).
\]
The first sum is $\binom{d+2}{3}$.  Since
$\sum_r|\log(1-q^{-r})|<\infty$, the second sum is $O_q(d)$.
Combining this with \eqref{eq:V-base-mass} gives
\eqref{eq:V-total-mass}.
\end{proof}

\begin{theorem}[Cubic escape of mass]\label{thm:escape-V}
Let $q\geq2$ and let $S$ be a fixed finite set of primes.  Then
\begin{equation}\label{eq:escape-V}
 \sum_{\substack{p\notin S\\p\nmid q}}v_p(V_d(1,q))\log p
 =\binom{d+2}{3}\log q+O_{q,S}(d^2).
\end{equation}
In particular, if $S$ contains the primes dividing $q$, then
\begin{equation}\label{eq:fixed-half-barrier}
 \sum_{p\in S}v_p(V_d(1,q))\log p
 =\binom{d+1}{3}\log q+O_{q,S}(d^2)
 =\left(\frac12+o(1)\right)\log V_d(1,q).
\end{equation}
\end{theorem}

\begin{proof}
For an odd fixed prime $p\nmid q$, let $t=\ord_p(q)$ and
$c=v_p(q^t-1)$.  The local contribution to the cyclotomic part of
\eqref{eq:V-factorization} is exactly
\begin{equation}\label{eq:weighted-local-profile}
 \sum_{m=1}^{\lfloor d/t\rfloor}(d+1-tm)\bigl(c+v_p(m)\bigr),
\end{equation}
by the lifting-the-exponent formula.  This is $O_{p,q}(d^2)$.  The analogous
$2$-adic expression, obtained from \eqref{eq:local-two}, is also $O_q(d^2)$.
A finite sum over $S$ is therefore $O_{q,S}(d^2)$.  Subtracting it from the
total cyclotomic mass \eqref{eq:V-cyc-mass} proves
\eqref{eq:escape-V}; adding the exact base contribution proves
\eqref{eq:fixed-half-barrier}.
\end{proof}

\begin{corollary}[Increasing residue orders]\label{cor:increasing-orders}
Fix $T\geq1$.  Among the cyclotomic prime factors contributing to
$V_d(1,q)$, the primes with $\ord_p(q)\leq T$ contribute only $O_{q,T}(d^2)$.
Consequently,
\[
 \sum_{\substack{p\nmid q\\\ord_p(q)>T}}
   v_p(V_d(1,q))\log p
 =\binom{d+2}{3}\log q+O_{q,T}(d^2).
\]
Thus asymptotically all of the cyclotomic half occurs at primes whose residue
orders tend beyond every fixed cutoff.
\end{corollary}

\begin{proof}
The primes with $\ord_p(q)\leq T$ divide the fixed integer
$\prod_{r=1}^{T}(q^r-1)$, so they form a finite set.  Apply
\cref{thm:escape-V}.
\end{proof}

\subsection{Interpretation for determinant arguments}

The change-of-basis matrix from monomials
$1,T,\ldots,T^d$ to the regular basis
$B_0^{(a,q)},\ldots,B_d^{(a,q)}$ is triangular.  Its determinant is
\[
 \prod_{k=1}^{d}\frac1{F_k(a,q)}=\frac1{V_d(a,q)}.
\]
Therefore $V_d(a,q)$ is not a convenient overestimate: it is the exact
covolume of the canonical integer-valued interpolation lattice.  The
half-mass theorem has the following sharp methodological consequence.

\begin{criterion}[Fixed-prime barrier]\label{crit:fixed-prime-barrier}
Suppose an auxiliary determinant on a $q$-geometric interpolation lattice
requires, at leading order, the full covolume gain $\log V_d(1,q)$.  Any
valuation argument restricted to a fixed finite set of primes can recover at
most
\[
 \left(\frac12+o(1)\right)\log V_d(1,q)
\]
from the universal geometric denominator.  Closing the remaining half
requires primes whose residue orders and identities vary with $d$, or a
non-$p$-adic mechanism of the same cubic strength.
\end{criterion}

This does not say that every fixed-prime proof is impossible: a sufficiently
strong archimedean estimate might need less than the full covolume.  It does
say that simply sharpening constants in the same fixed-$2$/$3$ valuation
bookkeeping cannot manufacture the missing cyclotomic half.

\section{Moving-prime compatibility: a zero--one endgame}

The half-mass theorem identifies where the missing arithmetic lives.  We now
measure whether the real-power monoid map preserves any of it.

\subsection{The compatibility mass}

For integers $q,Q>1$, define
\begin{equation}\label{eq:compat-mass}
 \mathfrak C_d(q,Q)
 :=\sum_{r=1}^{d}(d+1-r)
       \log\gcd(q^r-1,Q^r-1)
\end{equation}
and the total cyclotomic reference mass
\begin{equation}\label{eq:reference-mass}
 \mathfrak M_d(q)
 :=\sum_{r=1}^{d}(d+1-r)\log(q^r-1).
\end{equation}
By \eqref{eq:V-cyc-mass},
\begin{equation}\label{eq:M-asymptotic}
 \mathfrak M_d(q)=\binom{d+2}{3}\log q+O_q(d).
\end{equation}
The ratio $\mathfrak C_d/\mathfrak M_d$ records the weighted fraction of the
moving cyclotomic mass of $q$ that is also visible to $Q$.

The relevant external input is the theorem of
Bugeaud--Corvaja--Zannier \cite{BCZ2003}: if $q$ and $Q$ are multiplicatively
independent integers greater than one, then for every $\varepsilon>0$,
\[
 \log\gcd(q^r-1,Q^r-1)\leq\varepsilon r
\]
for all sufficiently large $r$.

\begin{theorem}[Compatibility zero--one law for a prime base]
\label{thm:compat-zero-one}
Let $q$ be prime and $Q\in\Z_{>1}$.
\begin{enumerate}
 \item If $Q=q^m$ for some $m\in\N$, then
 \begin{equation}\label{eq:compat-full}
   \mathfrak C_d(q,Q)=\mathfrak M_d(q)
 \end{equation}
 for every $d$.
 \item If $Q$ is not an integral power of $q$, then
 \begin{equation}\label{eq:compat-zero}
   \mathfrak C_d(q,Q)=o(d^3),
   \qquad
   \frac{\mathfrak C_d(q,Q)}{\mathfrak M_d(q)}\longrightarrow0.
 \end{equation}
\end{enumerate}
Thus the normalized compatibility mass tends to either $1$ or $0$; there is
no intermediate positive limiting proportion.
\end{theorem}

\begin{proof}
If $Q=q^m$, then $q^r-1$ divides $q^{mr}-1$, giving
\eqref{eq:compat-full}.  Otherwise $q$ and $Q$ are multiplicatively
independent.  Indeed, a relation $q^u=Q^v$ and unique factorization would
force $Q$ to be a power of the prime $q$.

Fix $\varepsilon>0$.  By the Bugeaud--Corvaja--Zannier theorem there is $R$
such that, for $r\geq R$,
\[
 \log\gcd(q^r-1,Q^r-1)\leq\varepsilon r.
\]
Hence
\begin{align*}
 \mathfrak C_d(q,Q)
 &\leq O_{q,Q,R}(d)
   +\varepsilon\sum_{r=R}^{d}(d+1-r)r\\
 &\leq O_{q,Q,R}(d)+\varepsilon\binom{d+2}{3}.
\end{align*}
After division by $d^3$, the limsup is at most $\varepsilon/6$.  Since
$\varepsilon$ is arbitrary, \eqref{eq:compat-zero} follows; the ratio statement
uses \eqref{eq:M-asymptotic}.
\end{proof}

\begin{corollary}[A concrete sufficient condition for \AE]
\label{cor:compat-AE}
Let $M=2^x$ and $N=3^x$ be integers.  If
\begin{equation}\label{eq:combined-positive-mass}
 \limsup_{d\to\infty}
 \frac{\mathfrak C_d(2,M)+\mathfrak C_d(3,N)}{d^3}>0,
\end{equation}
then $x\in\Z$.
\end{corollary}

\begin{proof}
Condition \eqref{eq:combined-positive-mass} implies that at least one of the
two compatibility masses is not $o(d^3)$.  By
\cref{thm:compat-zero-one}, either $M=2^m$ or $N=3^m$ for some
$m\in\N$.  In the first case $x=\log M/\log2=m$; in the second,
$x=\log N/\log3=m$.
\end{proof}

\begin{corollary}[Single-layer version]\label{cor:single-layer}
The same conclusion holds if, for one of the pairs $(q,Q)=(2,M)$ or $(3,N)$,
there is a constant $c>0$ and infinitely many $r$ such that
\[
 \log\gcd(q^r-1,Q^r-1)\geq c r.
\]
\end{corollary}

The support theorem of Corrales--Rodr\'{i}guez and Schoof
\cite{CorralesSchoof1997} gives a related qualitative endpoint: sufficiently
systematic transfer of prime supports or reduction orders from $q$ to $Q$
forces multiplicative dependence.  The weighted functional
\eqref{eq:compat-mass} is adapted more directly to the geometric Vandermonde
factorization, because the same weight $d+1-r$ occurs in
\eqref{eq:V-factorization}.

\subsection{The precise bridge conjecture}

\begin{conjecture}[Two-base moving-prime bridge]\label{conj:MP-bridge}
Let $M,N\in\Z_{>1}$ satisfy
\[
 \frac{\log M}{\log2}=\frac{\log N}{\log3}.
\]
Then
\[
 \limsup_{d\to\infty}
 \frac{\mathfrak C_d(2,M)+\mathfrak C_d(3,N)}{d^3}>0.
\]
\end{conjecture}

\begin{corollary}
\Cref{conj:MP-bridge} implies the \AE{} conjecture.
\end{corollary}

\begin{proof}
Apply \cref{cor:compat-AE}.
\end{proof}

The bridge conjecture is deliberately narrow.  It does not request full
support transfer, equality of $p$-adic exponents, or continuity of the power
map in every local topology.  A positive cubic proportion in either pair is
enough, and the zero--one theorem then upgrades that proportion to exact
multiplicative dependence.

\subsection{A determinant amplifier for gcd compatibility}

The next lemma explains why the gcd in \eqref{eq:compat-mass} is the correct
local object for determinants built from the multiplicative semigroup
$\langle q,Q\rangle$.

\begin{lemma}[Block-Vandermonde amplification]\label{lem:block-V}
Let $q,Q>1$, let $r\geq1$, and put
\[
 g_r=\gcd(q^r-1,Q^r-1).
\]
Choose distinct numbers
\[
 u_j=q^{a_j}Q^{b_j},\qquad a_j,b_j\in\N_0,
 \qquad 0\leq j\leq L.
\]
For any $n_0\in\N_0$, define
\[
 D=\det\bigl(u_j^{\,n_0+ir}\bigr)_{0\leq i,j\leq L}.
\]
Then
\begin{equation}\label{eq:block-divisibility}
 g_r^{\binom{L+1}{2}}\mid D.
\end{equation}
More precisely,
\begin{equation}\label{eq:block-factorization}
 D=\left(\prod_{j=0}^{L}u_j^{n_0}\right)
   \prod_{0\leq i<j\leq L}(u_j^r-u_i^r).
\end{equation}
\end{lemma}

\begin{proof}
Equation \eqref{eq:block-factorization} is the ordinary Vandermonde
factorization applied to $u_0^r,\ldots,u_L^r$.  Modulo $g_r$ one has
$q^r\equiv Q^r\equiv1$, hence $u_j^r\equiv1$ for every $j$.  Therefore every
pairwise difference in \eqref{eq:block-factorization} is divisible by $g_r$.
There are $\binom{L+1}{2}$ such differences.
\end{proof}

There is an equivalent finite-difference proof: on the arithmetic progression
$n_0,n_0+r,\ldots,n_0+Lr$, each forward difference of a sequence
$q^{an}Q^{bn}$ gains a factor $g_r$, and the $i$th transformed row gains
$g_r^i$.  The determinant amplification exponent is therefore
$0+1+\cdots+L$.

For the \AE{} data, one may use the same exponent pairs $(a_j,b_j)$ in the
two semigroups
\[
 u_j^{(2)}=2^{a_j}M^{b_j}=2^{a_j+b_jx},
 \qquad
 u_j^{(3)}=3^{a_j}N^{b_j}=3^{a_j+b_jx}.
\]
The real slopes $a_j+b_jx$ are identical while the local compatibility
factors are
\[
 \gcd(2^r-1,M^r-1),\qquad
 \gcd(3^r-1,N^r-1).
\]
This is the point at which the real common-exponent relation and the moving
prime factors meet in one determinant architecture.

\subsection{What remains to be proved in this route}

A complete moving-prime proof would have to reverse the usual direction of a
gcd estimate.  Known theorems give strong \emph{upper} bounds when the pairs
are independent.  The missing bridge must produce a \emph{lower} bound from
the simultaneous real relation.  Three possible concrete formulations are:
\begin{enumerate}
 \item prove \cref{conj:MP-bridge} directly;
 \item produce infinitely many $r$ satisfying the single-layer condition of
       \cref{cor:single-layer};
 \item construct a nonzero integer determinant whose divisibility lower bound
       contains a positive cubic portion of
       $\mathfrak C_d(2,M)+\mathfrak C_d(3,N)$, while the common real slope
       makes its archimedean size smaller than that divisor.
\end{enumerate}
The half-mass theorem shows that the third formulation has enough available
prime weight in principle.  It also shows that a construction that ignores
the varying factors $q^r-1$ cannot reach the required scale.

\section{A symmetric Kummer-period and coaction route}

The moving-prime program is unconditional but lacks the bridge lower bound.
The period route supplies a different kind of progress: it identifies the
exact standard conjectural statement that would make the logarithmic
determinant formal, after which unique factorization ends the proof.

\subsection{The symmetric prime symbol}

Let $S$ be a finite set of primes containing all prime divisors of $6MN$, and
let
\[
 E_S:=\bigoplus_{p\in S}\Q e_p.
\]
For a positive $S$-unit $A$, define its prime vector
\[
 \ell(A):=\sum_{p\in S}v_p(A)e_p.
\]
Let $u\odot v=(u\otimes v+v\otimes u)/2$ and consider
$\Sym^2(E_S)$.  Define the numerical quadratic-log map
\begin{equation}\label{eq:period-map-symbol}
 \pi_S:\Sym^2(E_S)\longrightarrow\R,
 \qquad
 \pi_S(e_p\odot e_q)=\log p\,\log q.
\end{equation}
The relevant symbol is
\begin{equation}\label{eq:Xi}
 \Xi(M,N):=e_2\odot\ell(N)-e_3\odot\ell(M).
\end{equation}
Then
\begin{equation}\label{eq:piXi}
 \pi_S(\Xi(M,N))
 =\log2\,\log N-\log3\,\log M.
\end{equation}

\begin{theorem}[Formal prime normal form]\label{thm:formal-normal-form}
For positive integers $M,N$,
\[
 \Xi(M,N)=0
\]
if and only if there is $k\in\N_0$ such that
\[
 M=2^k,\qquad N=3^k.
\]
\end{theorem}

\begin{proof}
Compare coefficients in the standard basis of $\Sym^2(E_S)$.  For every
$p\notin\{2,3\}$, the coefficient of $e_2\odot e_p$ is $v_p(N)$, so
$v_p(N)=0$; the coefficient of $e_3\odot e_p$ is $-v_p(M)$, so
$v_p(M)=0$.  The coefficient of $e_2\odot e_2$ is $v_2(N)$, and the
coefficient of $e_3\odot e_3$ is $-v_3(M)$, so these also vanish.  Finally,
the coefficient of $e_2\odot e_3$ is
\[
 v_3(N)-v_2(M),
\]
which must be zero.  Thus $M=2^k$ and $N=3^k$ with a common $k$.  The converse
is immediate.
\end{proof}

\begin{remark}[Why the symmetric square is essential]
The raw tensor
$e_2\otimes\ell(N)-e_3\otimes\ell(M)$ does not vanish for the genuine
solutions $M=2^k$, $N=3^k$: it becomes
$k(e_2\otimes e_3-e_3\otimes e_2)$.  Numerical multiplication of logarithms
kills this antisymmetric tensor because real multiplication is commutative.
Passing to $\Sym^2(E_S)$ removes exactly this unavoidable kernel and no more.
\end{remark}

\subsection{Restricted period injectivity}

\begin{definition}[Quadratic prime-log injectivity]\label{def:QPLI}
For a finite prime set $S$, let $\mathrm{QPLI}(S)$ denote the assertion that
the map $\pi_S$ in \eqref{eq:period-map-symbol} is injective.  Equivalently,
the products
\[
 \{\log p\,\log q:p,q\in S,\ p\leq q\}
\]
are linearly independent over $\Q$.
\end{definition}

\begin{theorem}[Conditional \AE{} theorem]\label{thm:conditional-QPLI}
Let $M=2^x$ and $N=3^x$ be positive integers, and let $S$ contain the prime
divisors of $6MN$.  If $\mathrm{QPLI}(S)$ holds, then $x\in\N_0$.
\end{theorem}

\begin{proof}
Equation \eqref{eq:logdet} and \eqref{eq:piXi} give
$\pi_S(\Xi(M,N))=0$.  Injectivity gives $\Xi(M,N)=0$.
By \cref{thm:formal-normal-form}, $M=2^k$ and $N=3^k$ for some
$k\in\N_0$.  Therefore $x=\log M/\log2=k$.
\end{proof}

Schanuel's conjecture implies $\mathrm{QPLI}(S)$: the numbers $\log p$ are
linearly independent over $\Q$ by unique factorization, and Schanuel then
predicts their algebraic independence because their exponentials are
algebraic.  The point of the present formulation is not to replace one open
problem by a larger one, but to isolate the exact weight-two subspace and its
formal coefficient extraction.

\subsection{Kummer periods and the weight-$(1,1)$ coaction}

For $A\in\Q_{>0}^{\times}$, the logarithm $\log A$ is a period of a Kummer
extension.  In the category of mixed Tate motives over
$\Z[S^{-1}]$, write schematically
\[
 L_A^{\mathfrak m}=\log^{\mathfrak m}(A),
 \qquad
 L_A^{\mathrm{dR}}=\log^{\mathrm{dR}}(A)
\]
for the motivic and de Rham logarithm classes.  The basic coaction is
primitive:
\begin{equation}\label{eq:primitive-coaction}
 \Delta L_A^{\mathfrak m}
 =L_A^{\mathfrak m}\otimes1+1\otimes L_A^{\mathrm{dR}}.
\end{equation}
See Deligne--Goncharov \cite{DeligneGoncharov2005} and Brown
\cite{Brown2017} for the mixed Tate and motivic-period frameworks.
Multiplicativity of the coaction gives the weight-$(1,1)$ component
\begin{equation}\label{eq:product-coaction}
 \Delta_{1,1}\bigl(L_A^{\mathfrak m}L_B^{\mathfrak m}\bigr)
 =L_A^{\mathfrak m}\otimes L_B^{\mathrm{dR}}
  +L_B^{\mathfrak m}\otimes L_A^{\mathrm{dR}}.
\end{equation}

Define the motivic determinant period
\begin{equation}\label{eq:motivic-D}
 \mathcal D^{\mathfrak m}(M,N)
 :=L_2^{\mathfrak m}L_N^{\mathfrak m}
   -L_3^{\mathfrak m}L_M^{\mathfrak m}.
\end{equation}
Its numerical period is exactly \eqref{eq:logdet}.  Using
$L_A=\sum_pv_p(A)L_p$, the symmetrization of
$\Delta_{1,1}\mathcal D^{\mathfrak m}$ is the prime symbol
$2\Xi(M,N)$.

\begin{theorem}[Motivic conditional proof]\label{thm:motivic-conditional}
Let $S$ contain the prime divisors of $6MN$.  Assume that the numerical
period map is injective on the weight-at-most-two motivic-period subspace of
the mixed Tate category generated by the Kummer classes
$\{L_p^{\mathfrak m}:p\in S\}$.  Then any real $x$ with
$2^x=M\in\Z$ and $3^x=N\in\Z$ is an integer.
\end{theorem}

\begin{proof}
The numerical period of $\mathcal D^{\mathfrak m}(M,N)$ is zero, so the
assumed injectivity gives
$\mathcal D^{\mathfrak m}(M,N)=0$.  Applying the weight-$(1,1)$ coaction and
using \eqref{eq:product-coaction}, then identifying the independent prime
Kummer classes in the motivic and de Rham factors, gives
$\Xi(M,N)=0$.  The conclusion follows from
\cref{thm:formal-normal-form}.
\end{proof}

The hypothesis is a restricted instance of the expected injectivity of the
period map in the Grothendieck--Kontsevich--Zagier philosophy
\cite{KontsevichZagier2001}.  The new structural point is the explicit
coaction computation: the project-level formal prime symbol is not merely an
analogy.  It is the weight-$(1,1)$ coaction shadow of the tensor-square Kummer
period whose numerical value is the logarithmic determinant.

\subsection{What the coaction does and does not solve}

The coaction extracts prime coefficients only after a numerical relation has
been lifted to a motivic relation.  That lifting is the unproved step.  It
would be circular to declare numerical vanishing sufficient without a period
injectivity theorem.  Nevertheless, the route gives three useful pieces of
information:
\begin{enumerate}
 \item it identifies the correct quotient (the symmetric square) and removes
       the false antisymmetric obstruction;
 \item it reduces every genuine solution to a coefficient comparison in
       weight one;
 \item it suggests that a moving-prime determinant should be viewed as an
       unconditional arithmetic surrogate for the coefficient-extracting
       coaction.
\end{enumerate}
The last point links the two main routes: the coaction sees all prime symbols
at once, while the half-mass theorem says that an elementary determinant must
let the relevant primes move with the scale to emulate that global
coefficient extraction.

\section{A function-field divisor analogue and the constant-field obstruction}

In a function field, valuations provide the missing coefficient detector
unconditionally.  This makes the function-field analogue easy and shows
precisely what is lost when the bases specialize to constants.

\subsection{Divisor content}

Let $C$ be a smooth projective geometrically integral curve over a field $K$,
and let $K(C)$ be its function field.  For $F\in K(C)^{\times}$ define
\begin{equation}\label{eq:divisor-content}
 c(F):=\gcd\{|v_P(F)|:P\in C,\ v_P(F)\neq0\}.
\end{equation}
If $F$ is constant, the set is empty; below we apply the definition only to
nonconstant functions.

\begin{theorem}[Function-field exponent rigidity]\label{thm:function-field}
Let $F,G,H,J\in K(C)^{\times}$ with $F$ and $G$ nonconstant, and let
$x\in\R$.  Suppose that the equalities of real divisors
\begin{equation}\label{eq:divisor-exponent}
 \operatorname{div}(H)=x\operatorname{div}(F),
 \qquad
 \operatorname{div}(J)=x\operatorname{div}(G)
\end{equation}
hold.  Then $x\in\Q$.  If $x=a/b$ in lowest terms, then
\[
 b\mid c(F),\qquad b\mid c(G).
\]
In particular, if $\gcd(c(F),c(G))=1$, then $x\in\Z$.
\end{theorem}

\begin{proof}
Choose a point $P$ with $v_P(F)\neq0$.  From the first equality in
\eqref{eq:divisor-exponent},
\[
 x=\frac{v_P(H)}{v_P(F)}\in\Q.
\]
Write $x=a/b$ in lowest terms.  For every $P$,
$b\,v_P(H)=a\,v_P(F)$; since $\gcd(a,b)=1$, one has
$b\mid v_P(F)$.  Thus $b\mid c(F)$.  The same argument applies to $G$.
If the two contents are coprime, then $b=1$.
\end{proof}

The criterion is sharp: if $F=U^b$ and $H=U^a$, then
$\operatorname{div}(H)=(a/b)\operatorname{div}(F)$.  For the natural
deformation of the bases,
\[
 F(t)=t,\qquad G(t)=t+1
\]
on $\mathbb P^1$, both divisor contents are one, so any rational-function
lift of $t^x$ and $(t+1)^x$ would force $x\in\Z$.

\subsection{Why specialization to $2$ and $3$ loses the proof}

\begin{lemma}[Constant-field kernel]\label{lem:constant-kernel}
If $C$ is projective and geometrically integral, then
\[
 \operatorname{div}(F)=0\qquad\Longleftrightarrow\qquad F\in K^{\times}.
\]
\end{lemma}

\begin{proof}
A rational function with no zeros or poles is a global invertible regular
function on a proper integral curve, hence is constant.
\end{proof}

At the special value $t=2$, the functions $t$ and $t+1$ become the constants
$2$ and $3$.  Their divisors vanish, and the coefficient detector in
\cref{thm:function-field} disappears entirely.  A deformation proof must
therefore supply more than a family passing through the four numbers
$(2,3,M,N)$.  It must lift the exponent relation to nonconstant rational or
motivic objects while controlling divisor content.  The equalities at a
single constant fiber do not produce such a lift.

This is the same obstruction seen in the period route.  Divisors make prime
coefficients visible over a function field; motivic coactions are expected to
make prime Kummer coefficients visible over the constant field; moving-prime
Vandermonde factors attempt to recover part of that visibility by elementary
congruences.

\section{Synthesis: a focused proof program}

The new results suggest a much narrower research plan than another broad
search through unrelated transcendence techniques.

\subsection{Stage I: build determinants that spend the cyclotomic half}

The geometric factorial and Vandermonde formulas specify the available local
budget exactly.  A candidate determinant should satisfy all of the following.
\begin{enumerate}
 \item Its interpolation nodes are organized into $2$- and $3$-geometric
       blocks, or into blocks in the semigroups $\langle2,M\rangle$ and
       $\langle3,N\rangle$.
 \item Its row operations expose factors $q^r-1$ with the natural weight
       $d+1-r$, rather than using only the base-prime factors
       $q^{\binom{d+1}{3}}$.
 \item The nonpolynomial columns created by the exponent $x$ are paired so
       that their loss of congruence is measured by
       $\gcd(q^r-1,Q^r-1)$.
 \item The common slope
       $\log M/\log2=\log N/\log3$ cancels the leading archimedean growth
       between the two base blocks.
\end{enumerate}
The block amplifier in \cref{lem:block-V} is a basic module for such a
construction.

\subsection{Stage II: prove a positive compatibility lower bound}

The decisive target is no longer an unspecified ``extra $p$-adic gain.''  It
is one of the explicit inequalities
\begin{align}
 \mathfrak C_d(2,M)+\mathfrak C_d(3,N)&\gg d^3
 \quad\text{along a subsequence},\label{eq:target-cubic}\\
 \log\gcd(2^r-1,M^r-1)+\log\gcd(3^r-1,N^r-1)&\gg r
 \quad\text{for infinitely many }r.\label{eq:target-linear}
\end{align}
Either statement closes the conjecture by
\cref{thm:compat-zero-one,cor:single-layer}.  The weighted version
\eqref{eq:target-cubic} is naturally aligned with the exact Vandermonde
covolume; the single-layer version \eqref{eq:target-linear} is more rigid and
may be easier to connect to support-problem methods.

\subsection{Stage III: use the period route as a structural benchmark}

The Kummer coaction says what a perfect coefficient detector would output:
\[
 e_2\odot\ell(N)-e_3\odot\ell(M).
\]
A determinant construction can be evaluated against this benchmark.  If its
moving-prime divisibility only sees aggregate heights and not the separate
prime coefficients, it is unlikely to distinguish a true integer exponent
from an unrelated multiplicatively independent pair.  The desired determinant
should approximate coefficient extraction by stratifying primes according to
$\ord_p(q)$ and by comparing which strata survive under $q\mapsto Q$.

\subsection{A precise theorem stack}

A complete proof along the moving-prime route can be decomposed into the
following stack.
\begin{longtable}{@{}p{0.19\textwidth}p{0.72\textwidth}@{}}
\toprule
Layer & Statement \\
\midrule
\endhead
A & Exact simultaneous $p$-ordering and denominator lattice on every
    geometric ray (proved in \cref{thm:simultaneous-ordering,thm:regular-basis}).\\
B & Half-mass and escape-of-mass at quadratic and cubic scale
    (proved in \cref{thm:half-mass-factorial,thm:escape-factorial,thm:half-mass-V,thm:escape-V}).\\
C & Compatibility zero--one law from gcd estimates
    (proved in \cref{thm:compat-zero-one}).\\
D & Determinant amplification of a given compatibility factor
    (proved in \cref{lem:block-V}).\\
E & Common-real-slope $\Rightarrow$ positive moving-prime compatibility
    (open bridge, \cref{conj:MP-bridge}).\\
F & Positive compatibility $\Rightarrow x\in\Z$
    (proved in \cref{cor:compat-AE}).\\
\bottomrule
\end{longtable}
Only Layer~E is missing from this stack.

\section{Sanity checks and small exact examples}

The exact formulas are simple enough to audit by hand.

\subsection{Generalized factorials}

For $q=2$ and $a=1$,
\[
 F_1=1,\qquad F_2=6,\qquad F_3=168,\qquad F_4=20160.
\]
Indeed,
\[
 F_4=2^{\binom42}(2-1)(2^2-1)(2^3-1)(2^4-1)
     =2^6\cdot1\cdot3\cdot7\cdot15.
\]
For $q=3$,
\[
 F_1=2,\qquad F_2=48,\qquad F_3=11232.
\]

\subsection{Vandermonde products}

For the nodes $1,2,4,8$,
\[
 V_3(1,2)
 =(2-1)(4-1)(8-1)(4-2)(8-2)(8-4)=1008.
\]
The product-of-factorials identity gives the same value:
\[
 F_1F_2F_3=1\cdot6\cdot168=1008.
\]
The factorization \eqref{eq:V-factorization} gives
\[
 2^{\binom43}(2-1)^3(2^2-1)^2(2^3-1)
 =2^4\cdot1^3\cdot3^2\cdot7=1008.
\]
For the nodes $1,3,9$,
\[
 V_2(1,3)=2\cdot8\cdot6=96=F_1(1,3)F_2(1,3).
\]

\subsection{Newton coefficient check}

For $k=2$,
\[
 [1,q,q^2]T^x
 =\frac{(q^x-1)(q^x-q)}{q(q-1)(q^2-1)}.
\]
At $x=0$ or $x=1$ the numerator vanishes, and at $x=2$ it equals the
denominator, as required for the polynomial $T^2$.  If $q^x=Q$ is an integer
but not a power of $q$, the numerator is a nonzero integer and
\cref{prop:base-dichotomy} predicts a growing base-prime denominator as $k$
increases.

\subsection{Compatibility extremes}

If $Q=q^m$, then
\[
 \gcd(q^r-1,Q^r-1)=q^r-1
\]
for every $r$, so the compatibility mass is full.  At the opposite extreme,
for multiplicatively independent $q,Q$, the Bugeaud--Corvaja--Zannier theorem
forces the normalized weighted mass to zero.  This checks that the functional
\eqref{eq:compat-mass} sharply distinguishes the desired conclusion rather
than merely measuring generic common small primes.

\section{Lean formalization plan}

The unconditional core is well suited to formalization because its main
identities are finite products and valuations.  The motivic section should be
formalized only as an abstract conditional interface until a suitable motive
library exists.

\subsection{Suggested module decomposition}

\begin{longtable}{@{}>{\raggedright\arraybackslash}p{0.25\textwidth}
                       >{\raggedright\arraybackslash}p{0.67\textwidth}@{}}
\toprule
File & Content \\
\midrule
\endhead
\path{GeomOrdering/Basic.lean}
 & Definitions of $G(a,q)$, the difference product $F_k(a,q)$, and the exact
   factorization in \eqref{eq:geom-factorial}.\\
\path{GeomOrdering/Gaussian.lean}
 & Recursive definition or polynomial evaluation of Gaussian binomial
   coefficients; proof that the quotient in \eqref{eq:ratio-qbinom} is an
   integer.\\
\path{GeomOrdering/POrdering.lean}
 & Simultaneous $p$-ordering theorem and the regular Newton--Gaussian basis.\\
\path{GeomOrdering/Valuation.lean}
 & The three cases in \cref{thm:local-profile}, with odd-prime and $2$-adic
   lifting lemmas isolated.\\
\path{GeomOrdering/Mass.lean}
 & Finite-product logarithm identities, convergence of $C_q$, and quadratic
   and cubic asymptotics.\\
\path{GeomOrdering/Vandermonde.lean}
 & Formula \eqref{eq:V-factorization}, product of generalized factorials, and
   fixed-finite-set escape estimates.\\
\path{Compatibility/WeightedGCD.lean}
 & Definitions of $\mathfrak C_d$ and $\mathfrak M_d$, the formal reduction
   from the BCZ upper-bound interface to the zero--one law, and
   \cref{cor:compat-AE}.\\
\path{Compatibility/BlockVandermonde.lean}
 & Determinant factorization and divisibility in \cref{lem:block-V}.\\
\path{PrimeSymbol/Symmetric.lean}
 & Finitely supported prime-exponent vectors, symmetric tensors, and the
   coefficient proof of \cref{thm:formal-normal-form}.\\
\path{FunctionField/ExponentDivisor.lean}
 & Divisor content and \cref{thm:function-field}, using the valuation API for
   function fields.\\
\path{Motivic/AbstractBridge.lean}
 & An abstract graded commutative period algebra with a primitive coaction;
   theorem that period injectivity plus the coaction implies the prime-symbol
   relation.\\
\bottomrule
\end{longtable}

\subsection{Core theorem sketches}

The first theorem can be represented without introducing $p$-orderings
initially:
\begin{verbatim}
theorem geom_diff_prod_dvd
    (a q n k : N) (hq : 2 <= q) (hk : k <= n) :
    product j in range k, (a*q^k - a*q^j)
      | product j in range k, (a*q^n - a*q^j)
\end{verbatim}
After proving the quotient is the evaluated Gaussian binomial coefficient,
$p$-ordering minimality is a one-line valuation consequence.

The weighted gcd theorem should isolate the deep external theorem behind a
clean interface:
\begin{verbatim}
def BCZGcdBound (q Q : N) : Prop :=
  forall eps > 0, exists R,
    forall r >= R,
      log (gcd (q^r - 1) (Q^r - 1)) <= eps * r
\end{verbatim}
The formal proof that this interface implies
$\mathfrak C_d(q,Q)=o(d^3)$ is elementary finite-sum analysis.  In this way,
formalization can certify the entire reduction even before BCZ itself is
formalized.

\subsection{Dependency order}

A practical order is
\[
 \text{finite products}
 \longrightarrow q\text{-binomial divisibility}
 \longrightarrow p\text{-ordering}
 \longrightarrow\text{local valuations}
 \longrightarrow\text{mass laws}
 \longrightarrow\text{compatibility reduction}.
\]
The prime-symbol normal form is independent and can be completed in parallel.
The function-field theorem is also independent.  This separation makes it
possible to obtain a machine-checked package of all unconditional progress
without waiting for the open bridge conjecture.

\section{Failure modes and falsification checklist}

A research report on an open problem is useful only if it states where its
new routes can fail.

\subsection{Moving-prime route}

\begin{enumerate}
 \item \textbf{Availability is not transfer.}  The half-mass theorem proves
       that many moving primes divide geometric denominators.  It does not
       prove that the map $q^n\mapsto Q^n$ preserves any of them.
 \item \textbf{Determinants usually give upper bounds for gcds.}  The
       divisibility in \cref{lem:block-V} turns a known gcd into a large
       divisor.  A proof still needs the common real slope to force a
       determinant small enough, or otherwise force the gcd large enough.
 \item \textbf{Primitive divisors alone are insufficient.}  Zsigmondy's
       theorem \cite{Zsigmondy1892} supplies new prime support for most
       $q^r-1$, but a single unweighted prime per layer need not carry cubic
       logarithmic mass.  The weighted mass, not merely the count of new
       primes, must enter.
 \item \textbf{A fixed set cannot simulate motion.}  Adding any finite list
       of auxiliary primes changes only the $O(d^2)$ error in
       \eqref{eq:fixed-half-barrier}.
\end{enumerate}

\subsection{Period route}

\begin{enumerate}
 \item \textbf{Numerical zero is not motivic zero unconditionally.}  The
       coaction can be applied to a motivic relation, not directly to an
       unexplained real-number equality.
 \item \textbf{The raw tensor is the wrong object.}  It incorrectly rejects
       the genuine integer solutions; symmetrization is mandatory.
 \item \textbf{The restricted injectivity is still deep.}  It is a precise
       weight-two period statement, not an established theorem for arbitrary
       prime sets.
\end{enumerate}

\subsection{Deformation route}

\begin{enumerate}
 \item A family whose generic fiber has primitive divisors would solve the
       exponent problem, but the constant fiber $(2,3)$ has zero divisor.
 \item Interpolating finitely many values does not produce a rational or
       motivic lift of $t^x$.
 \item Any claimed specialization argument must explicitly control divisor
       content through the degeneration; otherwise it silently assumes the
       missing bridge.
\end{enumerate}

\section{Conclusions}

The report establishes the following unconditional progress.
\begin{enumerate}
 \item The natural ordering of every integral geometric progression is a
       simultaneous $p$-ordering, with exact generalized factorial
       \[
       F_k(a,q)=a^kq^{\binom{k}{2}}\prod_{r=1}^k(q^r-1).
       \]
 \item The corresponding Newton--Gaussian polynomials form a canonical
       integral basis for all rational polynomials integer-valued on the
       progression.
 \item Fixed base primes carry exactly one half of the leading denominator
       mass, both at the single-factorial scale $k^2$ and at the full
       Vandermonde scale $d^3$.
 \item The other half escapes every fixed finite prime set and occurs at
       residue orders tending beyond every fixed bound.
 \item The weighted compatibility mass
       $\mathfrak C_d(q,Q)$ obeys a zero--one law for prime $q$: full mass for
       an integral power and $o(d^3)$ for multiplicative independence.
 \item A positive combined compatibility mass for $(2,M)$ and $(3,N)$ is an
       unconditional sufficient condition for the \AE{} conclusion.
 \item Block Vandermonde determinants amplify each compatibility factor by a
       quadratic exponent, giving a concrete mechanism for spending moving
       prime mass.
 \item The logarithmic determinant is the period of a symmetric tensor of
       Kummer logarithms; its weight-$(1,1)$ coaction is exactly the formal
       prime symbol whose vanishing forces $M=2^k$, $N=3^k$.
 \item A function-field divisor theorem proves the analogous exponent
       rigidity and isolates the constant-field kernel that blocks naive
       specialization.
\end{enumerate}

No unconditional proof of \eqref{eq:AE} is claimed.  The main open task has,
however, been reduced to a sharply testable statement rather than a vague
request for stronger estimates: prove that a common real exponent forces a
positive amount of moving cyclotomic compatibility in at least one of the
pairs $(2,M)$ and $(3,N)$.  The exact $p$-ordering and mass formulas show that
the required arithmetic budget exists; the BCZ zero--one theorem shows that
any positive cubic fraction would immediately force the desired integral
power; and the Kummer coaction identifies the formal coefficient relation
that an ideal determinant should emulate.

\appendix

\section{Two useful exact identities}

\subsection{Weighted local valuation in closed form}

For an odd prime $p\nmid q$, let $t=\ord_p(q)$,
$c=v_p(q^t-1)$, and $m=\lfloor d/t\rfloor$.  From
\eqref{eq:V-factorization} and lifting the exponent,
\begin{align}
 v_p(V_d(1,q))
 &=\sum_{r=1}^{d}(d+1-r)v_p(q^r-1)\\
 &=\sum_{u=1}^{m}(d+1-tu)\bigl(c+v_p(u)\bigr).
 \label{eq:appendix-weighted-vp}
\end{align}
This is the exact finite formula behind the $O(d^2)$ estimate for each fixed
cyclotomic prime.  It also suggests a computational stratification by residue
order $t$ rather than by prime size.

\subsection{Asymptotic Newton fingerprint}

If $Q=q^x$ and $x\notin\N_0$, then
\begin{equation}\label{eq:qpoch-limit}
 (-1)^k q^{-\binom{k}{2}}
 \prod_{r=0}^{k-1}(Q-q^r)
 \longrightarrow
 \prod_{r=0}^{\infty}(1-q^{x-r}),
\end{equation}
where the infinite product converges to a nonzero real number.  Indeed,
\[
 Q-q^r=-q^r(1-q^{x-r}),
\]
and $\sum_r|q^{x-r}|<\infty$ in the tail.  If $x=m\in\N_0$, the product is
identically zero for $k>m$.  Thus the Newton coefficient sequence has a sharp
zero/nonzero dichotomy, although its nonzero asymptotic by itself does not
supply the missing arithmetic contradiction.

\section{Minimal experimental agenda}

The formulas in the report suggest computations that are diagnostic rather
than merely numerical.
\begin{enumerate}
 \item For candidate integer pairs $(M,N)$ with
       $\log M/\log2\approx\log N/\log3$, compute the finite compatibility
       profiles
       \[
       r\longmapsto
       \bigl(\log\gcd(2^r-1,M^r-1),
             \log\gcd(3^r-1,N^r-1)\bigr).
       \]
       Approximate pairs should display the BCZ-independent regime; any
       anomalous mass concentration identifies the residue-order layers a
       determinant must exploit.
 \item Group the prime factors of $q^r-1$ by $\ord_p(q)=r$ and compare their
       survival under $Q^r-1$.  The weighted quantity in
       \eqref{eq:compat-mass} should be tracked, not only the number of common
       primes.
 \item Search over lattice exponent pairs $(a_j,b_j)$ in
       \cref{lem:block-V}.  Optimize the ratio between the guaranteed divisor
       $g_r^{\binom{L+1}{2}}$ and the exact archimedean determinant size under
       the common slope $a_j+b_jx$.
 \item Test multiscale blocks with several step sizes $r$.  Because valuations
       are assessed prime by prime, different residue orders can use different
       row orderings; a successful construction need not expose all moving
       primes in one literal sequence of row operations.
\end{enumerate}

\begin{thebibliography}{99}

\bibitem{AlaogluErdos1944}
L.~Alaoglu and P.~Erd\H{o}s,
\newblock On highly composite and similar numbers,
\newblock \emph{Transactions of the American Mathematical Society} 56 (1944),
448--469.

\bibitem{Andrews1976}
G.~E. Andrews,
\newblock \emph{The Theory of Partitions},
\newblock Addison--Wesley, 1976; reprinted by Cambridge University Press, 1998.

\bibitem{Bhargava1997}
M.~Bhargava,
\newblock $P$-orderings and polynomial functions on arbitrary subsets of
Dedekind rings,
\newblock \emph{Journal f\"ur die reine und angewandte Mathematik} 490 (1997),
101--127.

\bibitem{Bhargava2000}
M.~Bhargava,
\newblock The factorial function and generalizations,
\newblock \emph{American Mathematical Monthly} 107 (2000), 783--799.

\bibitem{BCZ2003}
Y.~Bugeaud, P.~Corvaja, and U.~Zannier,
\newblock An upper bound for the G.C.D. of $a^n-1$ and $b^n-1$,
\newblock \emph{Mathematische Zeitschrift} 243 (2003), 79--84.

\bibitem{Brown2017}
F.~Brown,
\newblock Notes on motivic periods,
\newblock \emph{Communications in Number Theory and Physics} 11 (2017),
557--655.

\bibitem{CahenChabert1997}
P.-J. Cahen and J.-L. Chabert,
\newblock \emph{Integer-Valued Polynomials},
\newblock Mathematical Surveys and Monographs 48, American Mathematical
Society, 1997.

\bibitem{CorralesSchoof1997}
C.~Corrales-Rodr\'{i}guez and R.~Schoof,
\newblock The support problem and its elliptic analogue,
\newblock \emph{Journal of Number Theory} 64 (1997), 276--290.

\bibitem{DeligneGoncharov2005}
P.~Deligne and A.~B. Goncharov,
\newblock Groupes fondamentaux motiviques de Tate mixte,
\newblock \emph{Annales scientifiques de l'\'Ecole Normale Sup\'erieure}
38 (2005), 1--56.

\bibitem{KontsevichZagier2001}
M.~Kontsevich and D.~Zagier,
\newblock Periods,
\newblock in \emph{Mathematics Unlimited---2001 and Beyond}, Springer, 2001,
771--808.

\bibitem{Waldschmidt2000}
M.~Waldschmidt,
\newblock \emph{Diophantine Approximation on Linear Algebraic Groups},
\newblock Grundlehren der mathematischen Wissenschaften 326, Springer, 2000.

\bibitem{Zsigmondy1892}
K.~Zsigmondy,
\newblock Zur Theorie der Potenzreste,
\newblock \emph{Monatshefte f\"ur Mathematik und Physik} 3 (1892), 265--284.

\end{thebibliography}

\end{document}
